\documentclass[book.tex]{subfiles}

\begin{document}

\chapter{Density Functional Theory}

\label{chap:dft}
\noindent\rule{11cm}{0.4pt} \\
\textbf{Prerequisites:} \autoref{chap:molecular_hamiltonian}, \autoref{chap:quantum_variational}  \\
\textbf{Difficulty Level:} ** \\
\noindent\rule{11cm}{0.4pt}
\vspace{5mm}

%\textcolor{red}{TODO: Define the electron density formally\\}
%\textcolor{red}{TODO: Define the Kohn-Sham orbital formally\\}


Let us posit that the quantum mechanical description of a system is completely defined by its wave function 
\begin{align}
\psi &= \psi\left(r_1,\dotsc, r_N, R_1,\dotsc, R_K, t \right),
\end{align}

which is a function of time $t$, electron coordinates $\left\{ \ec \right\}$ and nuclear coordinates $\left\{ \nc \right\}$.
The wave function satisfies the time dependent Schr\"{o}dinger equation
\begin{align}
\hat{H} \psi = i \frac{\partial}{\partial t} \psi , 
\end{align}
where $\hat{H}$ is the Hamiltonian operator. In coordinate representation, the Hamiltonian is given by
\begin{align}
\hat{H} &=
-\frac{1}{2}\sum_{i=1}^{N} \nabla _i^2
-\sum_{k=1} ^{K} \frac{1}{2M_k} \nabla _k^2
-\sum_{i=1}^{N} \sum_{k=1} ^{K} \frac{Z_k}{\left|r_i -R_k\right|} \\
& \qquad 
+\frac{1}{2} \sum_{i\neq j}^N\frac{1}{\left|r_i -r_j\right|} 
+\frac{1}{2} \sum_{k\neq l}^K\frac{Z_k Z_l}{\left|R_k -R_l\right|}. 
\end{align}
The first term is the kinetic energy of the electrons, the second term is the kinetic energy of the nuclei, the third term is the interaction between the electrons and the nuclei, the fourth term is the electron electron interaction and the fifth term is the interaction between the nuclei.
As the nuclei are 3 to 4 orders of magnitude heavier than the electrons, we can apply the Born-Oppenheimer (BO) approximation,
 where it is assumed that the electrons respond instantaneously to the movements of
the nuclei. This allows a separation of the problem into the electronic structure problem and the nuclear motion. Focusing on stationary solutions to the electronic structure problem, the time independent electronic wave function $\Phi = \Phi \left( \ec \right)$ satisfies the Schr\"{o}dinger equation
\begin{equation}
E \Phi = \left\{ \hat{T}_e +\hat{V}_{en} + \hat{V}_{ee} \right\} \Phi , 
\end{equation}
where $\hat{T}_e$ is the kinetic energy operator for the electrons, $\hat{V}_{en}$ is the electron nucleus interaction and $\hat{V}_{ee}$ is the electron electron interaction.

In density functional theory (DFT), the search for the $N$-particle wavefunction is replaced by a search for the electron density $n(\vec{r})$. This transformation reduces the problem from a $3N$ to a 3-dimensional problem.  From a physical point of view, it is clear that the ground state is uniquely defined by the external potential and the number of electrons in the potential, since the external potential specifies the Hamiltonian.

In 1964, Hohenberg and Kohn proved that the external potential is also uniquely determined by the ground state density $n(\vec{r})$.\cite{hohenbergkohn:1964} As the external potential determines the Hamiltonian, the ground-state $N$-particle wavefunction and all observables of the ground state are functionals of the density $n(\vec{r})$. In particular, the ground state energy is a functional of the density $E[n(\vec{r})]$. 

Using the variational principle, it can be shown,\cite{kohn:nobel} that the ground state density $n(\vec{r})$ minimizes the energy, $E[n(\vec{r})]$. Following the constrained search approach proposed by Lieb and Levy,\cite{Lieb:icqc:24,levy:preva:26} the ground state energy may be obtained by minimizing the energy over all possible wavefunctions with the density $n(\vec{r})$, and then minimizing over all densities.
The ground state energy may then be written formally as 
\begin{align}
E_0=\min_{n(\vec{r})} E[n(\vec{r}))]&=\min_{n(\vec{r})} \left\{ F[n(\vec{r})]+\int v_{ext}(\vec{r}) n(\vec{r}) \textup{d}\vec{r} \right\} \label{eq:Eafn} \, ,
\end{align}
where $v_{ext}$ is the external potential and $F[n(\vec{r})]$ is a universal but unknown functional obtained by minimizing over all wavefunctions with density $n(\vec{r})$

\begin{align}
F[n(\vec{r})]=\min_{\psi \rightarrow n(\vec{r})} \left< \psi \right|\hat{T}_e + \hat{V}_{ee} \left| \psi\right>.
\end{align}

The energy functional may thus be expressed as
\begin{align}
E[n(\vec{r})]=F[n(\vec{r})]+ V_{ext}[n(\vec{r})],  \label{eq:HKfunctional}
\end{align}
where $V_{ext}[n(\vec{r})]=\int v_{ext}(\vec{r}) n(\vec{r}) \textup{d}\vec{r}$. The first term is independent of the external potential and thus the same for all systems of electrons interacting through the Coulomb potential. The second term is easily evaluated and depends on the system.
The Hohenberg-Kohn theorem only proves that the energy is a functional of the density. However, it does not prescribe a scheme for obtaining this functional.  A year later, Kohn and Sham developed an approximate scheme for obtaining the functional~\cite{kohnsham:1965} and this will be discussed in the next section.

\subsection{The Kohn Sham Equations}
The Kohn-Sham scheme is based on the observation that we can construct a fictitious system based on non-interacting electrons in an effective potential, $v_{\rm eff}$, that can match the density of the true system consisting of interacting electrons.  Then, by invoking the Hohenberg-Kohn theorem, given the densities of the two systems are the same, then the energy and other ground state properties are the same.  Through the construction of the effective potential, we need to ensure that the electrons in the Kohn-Sham scheme have the exact same density as the (real) interacting electrons

The total energy, in the Kohn-Sham scheme, is given by\cite{kohnsham:1965}
\begin{align}
E\left[n(\vec{r}) \right] &= T_s\left[n(\vec{r}) \right] +V_{H}\left[n(\vec{r}) \right] +E_{xc}\left[n(\vec{r}) \right] + V_{ext}\left[n(\vec{r}) \right]
\end{align}
where $T_s\left[n(\vec{r}) \right]$ is the kinetic energy of the non-interacting electron gas with density $n(\vec{r})$, and $V_{H}\left[n(\vec{r}) \right]$ is the Hartree energy, given by
\begin{align}
V_{H}\left[n(\vec{r})\right] &= \frac{1}{2} \int \int \frac{ n(\vec{r^{\prime}})n(\vec{r}) }{ \left|\vec{r}-\vec{r}^{\prime}\right| } \textup{d} \vec{r}^{\prime} \textup{d}\vec{r}.
\end{align}
There is an additional ingredient required in the Kohn-Sham scheme to match the densities of the non-interacting system to that of the interacting system.  This term is called the exchange correlation energy functional, $E_{xc}[n]$, which is introduced as a correction that contains electron exchange and correlation. The exchange correlation functional is not known exactly, and the strength of density functional theory relies on the accuracy of the approximations to the exchange correlation functional, $E_{xc}$.

Using the variational principle invoking the constraint that the number of electrons needs to be conserved, we can derive the Kohn-Sham (KS) equations.  The KS equations are given by~\cite{kohnsham:1965}
\begin{align}
\left[-\frac{1}{2} \nabla^2 + v_{\rm eff}\right]\phi_j(\vec{r})=\epsilon \phi_j(\vec{r}), \label{eq:KSequation}
\end{align}
where $\phi_j(\vec{r})$ represents the Kohn-Sham orbital and the effective potential is
\begin{align}
v_{\rm eff}(\vec{r})&=v_{ext}(\vec{r})+\int \frac{ n(\vec{r^{\prime}})}{ \left|\vec{r}-\vec{r}^{\prime}\right| } \textup{d} \vec{r}^{\prime}+v_{xc}(\vec{r}), \label{eq:KSpotential}
%\intertext{the density is}
%n(\vec{r})&=\sum_{j=1}^{N} \left| \phi_j(\vec{r}) \right|^2
\intertext{and the exchange correlation potential is the functional derivative of $E_{xc}$}
v_{xc}(\vec{r})&=\frac{\delta E_{xc}}{\delta n(\vec{r})},
\end{align}
%
\begin{align}
\left[-\frac{1}{2} \nabla^2 + v_{ext}(\vec{r})+\int \frac{ n(\vec{r^{\prime}})}{ \left|\vec{r}-\vec{r}^{\prime}\right| } \textup{d} \vec{r}^{\prime}+v_{xc}(\vec{r}) \right]\phi_j(\vec{r})=\epsilon_j \phi_j(\vec{r}),
\end{align}

\begin{align}
n(\vec{r}) = \sum_{j \in occupied} \left| \phi_j(\vec{r}) \right|^2.
\end{align}

The KS equations describe non-interacting electrons in a fictitious potential that gives the same density as the interacting electrons.  Since the effective potential $v_{\rm eff}$ depends on $n(\vec{r})$, the Kohn-Sham equations \eqref{eq:KSequation} and \eqref{eq:KSpotential} must be solved iteratively until a self-consistent density is obtained.
%The KS equations must be solved iteratively until selfconsistency is obtained.

\subsection{Exchange Correlation Functionals}\label{dft:xc}
In their original scheme, Kohn and Sham proposed a simple approximation to the exchange correlation functional (xc-functional).\cite{kohnsham:1965}  They assumed that the exchange correlation energy at density $n(\vec{r})$ is the same as that of a homogeneous electron gas with the same density.  This later became known as the local-density approximation (LDA).

The exchange correlation energy, within the LDA, is given by
\begin{align}
E_{xc}^{LDA}=\int\epsilon_{xc}^{hom}(n(\vec{r}))n(\vec{r})\textup{d}\vec{r},
\end{align}
where $\epsilon_{xc}^{hom}(n(\vec{r}))$ is the energy of the exchange correlation hole in the homogeneous electron gas of density $n$.  The local density approximation is, in principle, unique.  %However, different parametrizations of $\epsilon_{xc}^{hom}$ have been proposed.\cite{ceperly:1980,wilk:1980,perdew:1981}
Despite its simplicity,  LDA has had remarkable success in describing properties such as lattice constants, vibrational frequencies and equilibrium geometries of physical systems. Dissociation energies of molecules and cohesive energies of solids are predicted within 10-20$\%$  and bond geometries are often predicted within 1$\%$.  The success of LDA is partly rationalized by the fact that the exchange correlation hole is normalized as in real systems.  It is exact in the limit of high density and the slowly varying limit.

Although LDA was often sufficiently accurate to describe properties of solids, the wide-spread application of density functional theory did not become commonplace until the development of the Generalized Gradient Approximations (GGA).  In a Generalized Gradient Approximation, the exchange correlation energy depends on the gradient of the density as well as the density.  The exchange correlation energy is given by 
\begin{align}
E_{xc}^{GGA}\left[ n(\vec{r})\right] &= \int \epsilon_{xc}^{GGA}\left( n(\vec{r}),\nabla n(\vec{r}) \right) \textup{d}\vec{r}.
\end{align}
However, unlike $\epsilon_{xc}^{hom}$, which is uniquely defined, $\epsilon_{xc}^{GGA}$ is not unique.  %Several GGA's have been suggested, for example PW91\cite{wang:prb:1991}, PBE\cite{ernzerhof:1999,ggamadesimple:1996} and RPBE\cite{hammer:prb:1999}. 
The choice of the functional depends on the system and properties of interest.

However, both LDA and GGA have serious shortcomings and we present a brief review of its failures to motivate the search for improvements.  A brief catalog of typical failures include:\cite{Albers:2009jq}
\begin{enumerate}
\item Band structure predictions of metallic materials that are experimentally known to be insulating (\textit{e.g.}, CoO and FeO), 
\item Absence of magnetism for materials that are magnetic (e.g., for many undoped high-temperature superconducting oxides) and vice versa (such as Pu), 
\item Band gaps that are much too small compared with experiment (\textit{e.g.}, for many semiconductors), 
\end{enumerate}
%More examples can be found elsewhere,\cite{Albers:2009jq} but this list suffices for our purposes.
When examining this list, it becomes clear that most of the problems listed are associated with the spectral properties of the electronic structure of a material. However, as explained carefully in the original papers on LDA,\cite{hohenbergkohn:1964,kohnsham:1965} this type of theory is designed to minimize the total ground-state energy of the electrons in a material as a functional of the spatial distribution of the number density of electrons. Thus, the eigenvalues of the Kohn-Sham equations\cite{kohnsham:1965} were never supposed to represent the actual quasi-particle spectrum of electrons. 

Nonetheless, because the eigenvalues often, in fact, are a reasonably good representation of the spectral properties measured in experiments, this identification is commonly made in practice. Hence, although this has no justification, most attempts to improvements to density functional theory beyond the GGA level are actually attempts to make corrections to the eigenvalue spectra to bring it into better agreement with various spectroscopic measurements that probe the quasi-particle properties of the materials, such as optical and photoemission spectra. Even the metal versus insulator problem involves this issue, since this distinction depends upon knowing the quasi-particle spectral distribution as a function of energy.


%\section{Notes}
%\begin{itemize}
%\item Hellman-Feynman theorems.
%\item Sternheimer equation. Connection with deep equilibrium models?
%\item Primal pass, solve SCF, f-rule, solve Sternheimer.
%\end{itemize}
% \subsection*{Methods}
% Using a standard force balance, the instantaneous power requirement for an electric vehicle in motion is given by Eq. \ref{eqn:power}.
% \begin{equation}
% \label{eqn:power}
% P =(\frac{1}{2} \rho C_D A v^3+C_{rr} mgv+t_f mgvZ)\cdot  \frac{1}{\eta_{bw}} +ma^+ v\cdot \frac{1}{\eta_{bw}} +ma^- v\cdot\eta_{bw} \eta_{brk}
% \end{equation}

% The equation includes aerodynamic, rolling friction, road grade, and inertial terms, including regenerative braking. $\eta_{bw}$ is the battery-wheels efficiency and $\eta_{brk}$ is the regenerative braking efficiency\cite{sripad_2017}. Of particular interest is the aerodynamic term,  $\frac{1}{2} \rho C_D A v^3$, which includes the non-dimensional drag coefficient ($C_D$) we compute using computational fluid dynamics (CFD). 



\end{document}